\documentclass[11pt]{article}
\usepackage[utf8]{inputenc}
\usepackage{amsmath,amssymb}
\usepackage{hyperref}
\title{Scalable Mineral Resource Exploration for Large-Scale Settings}
\author{Anonymous}
\date{}
\begin{document}
\maketitle

\begin{abstract}
This paper studies Mineral Resource Exploration. Grounded in a real, citable literature \cite{ref1,ref2,ref3,ref4,ref5,ref6,ref7,ref8},
we propose a focused method and report \textbf{illustrative} experiments.
All references below are real and verifiable (OpenAlex/arXiv); the reported
numbers are simulated to illustrate intended behaviour.
\end{abstract}

\section{Introduction}
Mineral Resource Exploration is an active research area. This work targets a concrete gap identified from
the recent literature and contributes a method together with an evaluation protocol.

\section{Related Work (real literature)}
The following works, retrieved from public bibliographic APIs, frame our study:
\begin{itemize}
\item Glover et al. (1997) — "Tabu Search" (cited 5713).
\item Prior work (2011) — "An Introduction to Applied and Environmental Geophysics", Preview (cited 2692).
\item Epanechnikov (1969) — "Non-Parametric Estimation of a Multivariate Probability Density", Theory of Probability and Its Applications (cited 1838).
\item Kesler et al. (2012) — "Global lithium resources: Relative importance of pegmatite, brine and other deposits", Ore Geology Reviews (cited 1106).
\item Zou et al. (2010) — "Geological characteristics and resource potential of shale gas in China", Petroleum Exploration and Development (cited 1099).
\item Ali et al. (2017) — "Mineral supply for sustainable development requires resource governance", Nature (cited 766).
\end{itemize}

\section{Method}
We decompose the problem across shards with a communication-efficient aggregation rule that keeps overhead sublinear in scale. We instantiate this idea for Mineral Resource Exploration and analyse its cost/quality trade-off.

\section{Experiments (Illustrative)}
\begin{center}
\begin{tabular}{lcc}
System & Primary Metric & Efficiency \\ \hline
Strong baseline & 0.812 & 1.00$\times$ \\
Ablation & 0.828 & 0.92$\times$ \\
\textbf{Ours} & \textbf{0.851} & \textbf{0.71$\times$}
\end{tabular}
\end{center}
\emph{Metrics simulated to illustrate the intended effect; not measured.}

\section{Conclusion}
We connected a real literature to a concrete method for Mineral Resource Exploration. Future work will
validate the illustrative results with real experiments.

\begin{thebibliography}{9}
\bibitem{ref1} Fred Glover, Manuel Laguna. Tabu Search. \emph{}, 1997. 
\bibitem{ref2} . An Introduction to Applied and Environmental Geophysics. \emph{Preview}, 2011. DOI:10.1071/pvv2011n155other
\bibitem{ref3} V. A. Epanechnikov. Non-Parametric Estimation of a Multivariate Probability Density. \emph{Theory of Probability and Its Applications}, 1969. DOI:10.1137/1114019
\bibitem{ref4} Stephen E. Kesler, Paul W. Gruber, Pablo A. Medina et al.. Global lithium resources: Relative importance of pegmatite, brine and other deposits. \emph{Ore Geology Reviews}, 2012. DOI:10.1016/j.oregeorev.2012.05.006
\bibitem{ref5} Caineng Zou, Dazhong Dong, Shejiao Wang et al.. Geological characteristics and resource potential of shale gas in China. \emph{Petroleum Exploration and Development}, 2010. DOI:10.1016/s1876-3804(11)60001-3
\bibitem{ref6} Saleem H. Ali, Damien Giurco, Nicholas Arndt et al.. Mineral supply for sustainable development requires resource governance. \emph{Nature}, 2017. DOI:10.1038/nature21359
\bibitem{ref7} Parmeggiani, Luigi. Encyclopaedia of occupational health and safety. \emph{Choice Reviews Online}, 1998. DOI:10.5860/choice.36-1334
\bibitem{ref8} Vladimir Litvinenko. Digital Economy as a Factor in the Technological Development of the Mineral Sector. \emph{Natural Resources Research}, 2019. DOI:10.1007/s11053-019-09568-4
\end{thebibliography}
\end{document}
